\documentclass[12pt,a4paper]{article}
\usepackage{a4wide}
\usepackage{graphicx}
\usepackage{float}
\usepackage{booktabs,caption}
\usepackage[flushleft]{threeparttable}
\usepackage{tabularx}
\usepackage[nottoc]{tocbibind} 
\usepackage{amssymb,amsmath}
\usepackage{pifont}% http://ctan.org/pkg/pifont
\usepackage[margin=1in]{geometry}
\usepackage{hyperref}
\usepackage{graphicx,grffile}
\usepackage{longtable} 
\usepackage{multirow}
\usepackage{wrapfig}
\usepackage{colortbl}
\usepackage{pdflscape}
\usepackage{tabu}

\newcommand{\cmark}{\ding{51}}
\newcommand{\xmark}{\ding{55}}

\hypersetup{unicode=true, pdftitle={NCLT Cause Lists Codebook},
  pdfauthor={Ami Dagli and Rahul Somani}, pdfborder={0 0 0}, breaklinks=true}
\urlstyle{same} % don't use monospace font for urls

% Commands for the FRG note look -- Do Not Change, Do Not Add -- %

\makeatletter
\def\maxwidth{\ifdim\Gin@nat@width>\linewidth\linewidth\else\Gin@nat@width\fi}
\def\maxheight{\ifdim\Gin@nat@height>\textheight\textheight\else\Gin@nat@height\fi}
\makeatother
\setkeys{Gin}{width=\maxwidth,height=\maxheight,keepaspectratio}
\IfFileExists{parskip.sty}{%
  \usepackage{parskip} }{% else
  \setlength{\parindent}{0pt} \setlength{\parskip}{6pt plus 2pt minus
    1pt} } \setlength{\emergencystretch}{3em} % prevent overfull lines
\providecommand{\tightlist}{%
  \setlength{\itemsep}{0pt}\setlength{\parskip}{0pt}}

\setcounter{secnumdepth}{5}
% Redefines (sub)paragraphs to behave more like sections
\ifx\paragraph\undefined\else \let\oldparagraph\paragraph
\renewcommand{\paragraph}[1]{\oldparagraph{#1}\mbox{}} \fi
\ifx\subparagraph\undefined\else \let\oldsubparagraph\subparagraph
\renewcommand{\subparagraph}[1]{\oldsubparagraph{#1}\mbox{}} \fi

\let\rmarkdownfootnote\footnote%
\def\footnote{\protect\rmarkdownfootnote}

\newlength{\outermargingap}
\setlength{\outermargingap}{1in} \newlength{\textwidthorig}
\setlength{\textwidthorig}{\textwidth}
\renewcommand{\baselinestretch}{1.1}

\usepackage{datetime}
\renewcommand{\dateseparator}{-} \newcommand{\todayiso}{\the\year
  \dateseparator \the\month \dateseparator \the\day}

% Commands for the title page of the FRG note look starts below
% Put in the title, authors and date here.

\begin{document}

\begin{titlepage}


  \vspace*{\stretch{1}}
  \parbox{\textwidthorig}{ \hrule % For horizontal lines
    \vspace{\baselineskip}
    \center{\Large Instruction manual for the FRG NCLT cause-list data-set}\\[1ex]
    July 2020 \vspace{\baselineskip} % Line spacing
    \hrule % For horizontal lines
  }

  \vspace{\stretch{1}}
  \vspace{\stretch{1}}
  
  \parbox{\textwidthorig}{ \center
    Finance Research Group\\[1ex]
    Mumbai}
\end{titlepage}

\newpage

\section{System requirements}

The ``FRG NCLT cause-list database system'' uses a set of Python and R programs. The python programs are first used to web-scrape data from NCLT court-bench websites to create a directory of cause-list files called \texttt{NCLT-DATA}. This is a time-consuming part of the system which is fraught with missteps and failures because of the lack of reliable access to the NCLT websites. This requires a dedicated system and dedicated manpower to follow-up when the access fails to repeat steps until the complete data is acquired

Second, the R programs then extract data to create an aggregated research-ready cause-list data-set in a file called \texttt{frg\_nclt\_cause-list\_dataset.csv}.  Once the source cause-list data is extracted from the NCLT websites in the first step, this is an easy, low-cost and reliable process.

The Python programs required include:
\begin{enumerate}
\item Python version 3.7 or higher.
\item Selenium Virtual Browser for either Geckodriver or Firefox.
\item System libraries including \texttt{libcurl4-openssl-dev, libxml2-dev, openssl, libssl-dev, libmagick++-dev}.
\end{enumerate}

The R programs required include:
\begin{enumerate}
\item R version 3.6.2 (or higher)
\item R packages:

  $\bullet$ tidyverse
  $\bullet$ httr
  $\bullet$ rvest
  $\bullet$ openssl
  $\bullet$ xml2
  $\bullet$ dplyr
  $\bullet$ glue
  $\bullet$ purrr
  $\bullet$ stringr
  $\bullet$ readr
  $\bullet$ optparse
  $\bullet$ ggplot2
  $\bullet$ tidyr
  $\bullet$ data.table
  $\bullet$ knitr
  $\bullet$ kableExtra
  $\bullet$ magick
  $\bullet$ rmarkdown
\end{enumerate}

\textbf{NOTE:} This system has been developed and tested for Linux (Debian 10).

\section{Running the system: Creating the directory of cause-list files}

The following commands should be run in the given sequence.

\begin{description}
\item[Step 1.] Update the file \texttt{benches\_courts\_names.pkl}
  which stores the list of the current NCLT courts and benches using a
  Python program: \texttt{\_01\_get\_benches\_courts.py}

  \begin{itemize}
  \item \textsc{Command}: \texttt{python3
      \_01\_get\_benches\_courts.py}
  \item \textsc{Output}: \texttt{benches\_courts\_names.pkl}
  \end{itemize}


\item[Step 2.] Validate the list in
  \texttt{benches\_courts\_names.pkl} using a Python program
  \texttt{\_02\_verify\_bench\_courts.py}

  \begin{itemize}
  \item \textsc{Command}: \texttt{python3
      \_02\_verify\_bench\_courts.py}
  \item \textsc{Output}: \texttt{benches\_courts-VALIDATED.pkl}
  \end{itemize}

  \textsc{Success Indicator:} If
  \texttt{benches\_courts-VALIDATED.pkl} is not created, re-do
  \textbf{Step 1} until it gets created. This may need to be done over
  a staggered period of time to adjust for problems in accessing the
  SOURCE location.

\item[Step 3.] Scrape URLs for the bench-court combinations listed in
  \texttt{benches\_courts-VALIDATED.pkl} using a Python program
  \texttt{\_03\_get\_bench\_court\_urls.py}.

  \begin{itemize}
  \item \textsc{Command}:\texttt{python3
      \_03\_get\_bench\_court\_urls.py}
  \item \textsc{Output}: \texttt{bench\_court\_urls.csv}, \texttt{bench\_court\_skipped.csv}
  \end{itemize}

  \textsc{Success indicator:}

  If the program is successful, \texttt{bench\_court\_urls.csv} is
  created.

  If the program is not successful, \texttt{bench\_court\_skipped.csv}
  is created.

  Repeat this step till \texttt{bench\_court\_urls.csv} is created.

\item[Step 4.] Extract cause-list data for a day from the list of
  bench-courts in \texttt{bench\_court\_urls.csv} using the Python
  program \texttt{\_04\_get\_cause\_lists.py}.

  The program uses \texttt{bench\_courts\_url.csv} as the default
  source for the bench-court list, and requires start date, end date
  and a path to save the extracted files as input parameters.

  \begin{itemize}
  \item \textsc{Command}:
  
    \texttt{python3 \_04\_get\_cause\_lists.py --start\_date YYYYMMDD
      --end\_date YYYYMMDD --data\_dir /directory/path/}
  
  \item \textsc{Output}: Sub-directories for each bench-court as
    listed in \texttt{bench\_courts\_url.csv}
    (e.g. chennai-bench\_court-i), which further contain
    subdirectories \texttt{data} and \texttt{logs}.
  \end{itemize}

  \textsc{Success indicator:}

  The \texttt{data} directory will have
  YYYY-MM-DD\_location-bench\_court-i.csv files for each date between
  \texttt{start-date} and \texttt{end-date}.

  The \texttt{log} directory will have files starting with
  \texttt{no\_data\_[bench-court-name].csv} and
  \texttt{verify\_dates\_[bench-court-name].csv} for missing files.

  If there are files created in \texttt{logs}, these can be used to
  manually follow up on missing cause-list data files.

\end{description}


\section{Running the system: Aggregating cause-list files into the research data-set}

A single data-set of cases listed in the cause-list files across
bench-courts and days is aggregated into a single file using two R
programs \texttt{01\_generate\_cpca\_numbers.R} and
\texttt{02\_generate\_dataset\_from\_cpca\_numbers.R}, which are
executed in sequence.


If the directory where the scraped NCLT cause-list files are stored is
\texttt{SOURCEPATH}, then the following commands should be run to
create the NCLT cause-list research data-set:\footnote{The command
  \texttt{Rscript 01\_generate\_cpca\_numbers.R -h} prints out the
  list of parameters that the user can set for
  \texttt{01\_generate\_cpca\_numbers.R -h}

  The command \texttt{Rscript
    02\_generate\_dataset\_from\_cpca\_numbers.R -h} prints out the
  list of parameters that can be set for
  \texttt{02\_generate\_dataset\_from\_cpca\_numbers.R} }
\begin{enumerate}
\item \texttt{Rscript 01\_generate\_cpca\_numbers.R --path\_data
    SOURCEPATH}
\item \texttt{Rscript 02\_generate\_dataset\_from\_cpca\_numbers.R
    --path\_data SOURCEPATH --input\_cpca\_file 01\_cpca\_numbers.txt}
\end{enumerate}

The first command creates the file \texttt{01\_cpca\_numbers.txt} as
output.

This file is taken as input by the second program to create
\texttt{frg\_nclt\_causelist\_dataset.csv} which is the final research
dataset.

The description of the data created and stored in
\texttt{frg\_nclt\_causelist\_dataset.csv} can be found in the
\texttt{Codebook} which is generated using the steps in the following
section.



\section{Generating the Codebook describing the research data-set}

In order to generate the Codebook which describes the research
data-set, you will need the following files:
\texttt{codebook\_doc.Rmd, codebook\_generator.R,
  codebook\_preprocessor.R}

The command \texttt{Rscript codebook\_generator.R --path\_data
  SOURCEPATH} will generate the file \texttt{./CODEBOOK/codebook.pdf}

\end{document}
